\documentclass[article,nojss]{jss}

%% --- Supplementary material for the topica manuscript. Same arXiv-first build
%% as topica.tex (the `nojss` option drops the journal masthead, so only jss.cls
%% and jss.bst are needed). For a JSS journal submission this is the supplementary
%% PDF that accompanies the main article; remove `nojss` alongside the main file.

\usepackage{amsmath,amssymb,booktabs,tabularx}

\author{Neal Caren\\University of North Carolina at Chapel Hill}
\Plainauthor{Neal Caren}

\title{Supplement to \pkg{topica}: Side by Side with the Reference
Implementations}
\Plaintitle{Supplement to topica: Side by Side with the Reference Implementations}
\Shorttitle{\pkg{topica} supplement: side-by-side validation}

\Abstract{
This supplement accompanies ``\pkg{topica}: A Unified \proglang{Python} Framework
for Topic Modeling in the Social Sciences.'' Section~5 of the main article
summarizes the cross-implementation agreement as a single alignment cosine per
model; here we show the underlying picture in full. For each reference package
\pkg{topica} validates against, both the original package and \pkg{topica} are fit
on the \emph{same} poliblog corpus with the same settings, and their output is
displayed twice, side by side, so a reader can compare the recovered topics word
for word rather than through a summary statistic. The tables are generated by
\code{paper/gen\_validation\_appendix.py} as part of \code{paper/reproduce.py}.
}

\Keywords{topic models, validation, reproducibility, \pkg{topica}}
\Plainkeywords{topic models, validation, reproducibility, topica}

\Address{
  Neal Caren\\
  Department of Sociology\\
  University of North Carolina at Chapel Hill\\
  E-mail: \email{neal.caren@unc.edu}\\
  URL: \url{https://nealcaren.github.io/topica/}
}

\begin{document}

\section{Side by side with the reference implementations}\label{app:sidebyside}

Section~5 of the main article summarizes the cross-implementation agreement as a
single alignment cosine per model. This supplement shows the underlying picture: for
each reference package \pkg{topica} validates against, both the original package
and \pkg{topica} are fit on the \emph{same} poliblog corpus with the same
settings, and their output is displayed twice, side by side. The fifteen legs span
the classical roster---LDA, NMF, LSA, STM, CTM, content/SAGE, DMR, g-DMR, keyATM,
supervised LDA, labeled LDA, HDP, Pachinko allocation, the dynamic topic model, and
the structural topic-and-sentiment model---each against an established
implementation: Java \pkg{MALLET} and \pkg{tomotopy} (LDA), \pkg{scikit-learn} (NMF,
LSA), \proglang{R}~\pkg{stm} (STM, CTM, content), \proglang{R}~\pkg{keyATM},
\pkg{tomotopy} (DMR, g-DMR, sLDA, labeled LDA, HDP, Pachinko, DTM), and the authors'
\proglang{R} reference (STS). Each subsection aligns the two engines' topics (by
topic-word cosine, except keyATM keyword topics and labeled-LDA label topics, which
are index-aligned by construction) and reports the top words from each, plus the
feature that distinguishes that model: covariate prevalence effects (STM, DMR),
continuous-time prevalence (g-DMR), topic correlations (CTM), per-group wording
(content/SAGE), seeded keywords (keyATM), a response regression (sLDA), label-topic
correspondence (labeled LDA), the inferred topic count (HDP), the super/sub-topic
hierarchy (Pachinko), word drift over time (DTM), and covariate-driven sentiment
(STS). LDA is shown against two independent references at once (Java \pkg{MALLET}
and \pkg{tomotopy}).

Agreement falls into three regimes. The spectral-initialized structural models (STM,
CTM, content) and the deterministic factorizations (NMF, LSA) start from the same
solution as their reference and coincide almost exactly (cosine $\geq 0.95$); the
supervised models constrained by labels likewise match near-exactly. For the
random-initialized collapsed-Gibbs models (LDA, DMR, g-DMR, sLDA, HDP, Pachinko,
DTM) the agreement is statistical, as Section~5 of the main article notes: the
dominant topics coincide while the diffuse, low-prevalence tail topics, which differ
between any two random-initialized runs, account for most of the alignment gap. For
those legs the relevant yardstick is not an absolute cosine but the reference's own
seed-to-seed agreement---the ceiling any two independent samplers can reach---which
the subsections report alongside the cross-implementation figure (DMR, for instance,
sits at the level two \pkg{tomotopy} runs reach with each other). For the
nonparametric HDP the headline result is instead that both engines infer nearly the
same topic count (about thirty) without being told it, since HDP topic identities are
seed-sensitive even within one engine.

The neural and embedding models (ETM, ProdLDA, CombinedTM, ZeroShotTM, InfoCTM,
FASTopic, the dynamic ETM, BERTopic, Top2Vec, the embedding LDA) and the LLM-based
TopicGPT are validated in Section~5 of the main article and \code{parity/} against
their own reference implementations on the corpora those references use, and so are
not part of this poliblog side-by-side. The tables here are generated by
\code{paper/gen\_validation\_appendix.py} as part of \code{reproduce.py}; legs whose
reference toolchain (\proglang{R} with \pkg{stm}/\pkg{keyATM}, the \pkg{MALLET}
CLI, \pkg{scikit-learn}, \pkg{tomotopy}, or the \pkg{STS} replication package) was
absent at build time are marked skipped.

% Generated by paper/gen_validation_appendix.py — do not edit by hand.

\noindent Table~\ref{tab:app:map} maps every \pkg{topica} model to its reference, corpus, and validation check. The classical models follow as side-by-side legs; the rest are validated on their reference's native corpus.

\begin{table}[p]\centering\footnotesize
\caption{Validation map: every \pkg{topica} model, its reference implementation, the corpus on which agreement is measured, the check applied, and where the detailed comparison lives. The classical models are the side-by-side legs of this appendix; the rest are validated on their reference's native corpus in \code{parity/} or \code{tests/} (those scripts compute the figures at run time). Random-initialized samplers are read against the reference's own seed-to-seed agreement, not an absolute threshold.}
\label{tab:app:map}
\begin{tabular}{@{}l >{\raggedright\arraybackslash}p{0.16\textwidth} >{\raggedright\arraybackslash}p{0.12\textwidth} >{\raggedright\arraybackslash}p{0.24\textwidth} >{\raggedright\arraybackslash}p{0.20\textwidth}@{}}
\toprule
model & reference & corpus & agreement & detail \\
\midrule
\addlinespace \multicolumn{5}{@{}l}{\textit{Classical (poliblog side-by-side, this appendix)}} \\[1pt]
LDA & \pkg{MALLET}, \pkg{tomotopy} & poliblog & 0.82 / 0.80 & Tab.~\ref{tab:app:lda} \\
NMF & \pkg{scikit-learn} & poliblog & 0.98 & Tab.~\ref{tab:app:nmf} \\
LSA & \pkg{scikit-learn} & poliblog & 0.99 & Tab.~\ref{tab:app:lsa} \\
STM & \proglang{R}~\pkg{stm} & poliblog & 0.99 & Tab.~\ref{tab:app:stm} \\
CTM & \proglang{R}~\pkg{stm} & poliblog & 0.99 & Tab.~\ref{tab:app:ctm} \\
SAGE/content & \proglang{R}~\pkg{stm} & poliblog & 0.96 & Tab.~\ref{tab:app:content} \\
DMR & \pkg{tomotopy} & poliblog & 0.82 (ceil.\ 0.71) & Tab.~\ref{tab:app:dmr} \\
g-DMR & \pkg{tomotopy} & poliblog & 0.82 & Tab.~\ref{tab:app:gdmr} \\
KeyATM & \proglang{R}~\pkg{keyATM} & poliblog & 0.89 & Tab.~\ref{tab:app:keyatm} \\
sLDA & \pkg{tomotopy} & poliblog & 0.73 & Tab.~\ref{tab:app:slda} \\
LabeledLDA & \pkg{tomotopy} & poliblog & 1.00 & Tab.~\ref{tab:app:labeledlda} \\
HDP & \pkg{tomotopy} & poliblog & count 37$\approx$28 & Tab.~\ref{tab:app:hdp} \\
PA & \pkg{tomotopy} & poliblog & 0.75 & Tab.~\ref{tab:app:pa} \\
DTM & \pkg{tomotopy} & poliblog & 0.73 & Tab.~\ref{tab:app:dtm} \\
STS & \proglang{R} (authors) & poliblog & 0.93 & Tab.~\ref{tab:app:sts} \\
\addlinespace \multicolumn{5}{@{}l}{\textit{Neural (autoencoding / optimal-transport)}} \\[1pt]
ProdLDA & PyTorch AVITM & 20\,Newsgroups & $c_v$, $c_{\text{npmi}}$, cross-NMI within reference seed noise & \code{parity/prodlda\_compare.py} \\
CombinedTM & PyTorch CTM & synthetic & $c_v$, $c_{\text{npmi}}$, cross-NMI & \code{parity/combinedtm\_compare.py} \\
ZeroShotTM & PyTorch CTM & synthetic + emb. & $c_v$, $c_{\text{npmi}}$, cross-NMI & \code{parity/zeroshot\_compare.py} \\
InfoCTM & PyTorch InfoCTM & bilingual synth. & cross-lingual alignment $\geq 0.8$ & \code{parity/infoctm\_compare.py} \\
FASTopic & \pkg{fastopic} & 20\,Newsgroups & $c_v$, $c_{\text{npmi}}$, diversity & \code{parity/fastopic\_compare.py} \\
ETM & PyTorch ETM & synthetic & finite-diff.\ gradient; coherence & \code{tests/test\_etm.py} \\
DETM & reference fit & synthetic time series & Hungarian topic cosine $\geq 0.7$ & \code{tests/test\_detm.py} \\
\addlinespace \multicolumn{5}{@{}l}{\textit{Embedding-based}} \\[1pt]
BERTopic & \pkg{BERTopic} & planted clusters & adj.\ Rand index, purity & \code{parity/top2vec\_compare.py} \\
Top2Vec & \pkg{BERTopic} & planted clusters & adj.\ Rand index, purity & \code{parity/top2vec\_compare.py} \\
EmbeddingLDA & structural & planted emb.\ blocks & block recovery & \code{tests/test\_embedding\_lda.py} \\
\addlinespace \multicolumn{5}{@{}l}{\textit{LLM-based}} \\[1pt]
TopicGPT & deterministic backend & synthetic & taxonomy + assignment & \code{tests/test\_topicgpt.py} \\
\addlinespace \multicolumn{5}{@{}l}{\textit{Short-text / seeded}} \\[1pt]
GSDMM & MGP (Yin \& Wang) & planted short texts & cluster recovery, inferred $K$ & \code{tests/test\_gsdmm.py} \\
PT & structural & synthetic & shape / contract & \code{tests/test\_extra\_models.py} \\
SeededLDA & \pkg{seededlda} & planted seeds & seed-word placement & \code{tests/test\_seeded\_warp.py} \\
\addlinespace \multicolumn{5}{@{}l}{\textit{Hierarchical / experimental}} \\[1pt]
HLDA & nested CRP (no poliblog ref.) & synthetic & structural (tree, paths) & \code{tests/} \\
ECTM & experimental tier; own paper & planted group$\times$time & content drift + placebo & \code{tests/test\_ectm.py} \\
\bottomrule
\end{tabular}
\end{table}

\subsection[Latent Dirichlet allocation (vs Java MALLET and tomotopy)]{Latent Dirichlet allocation (vs Java \pkg{MALLET} and \pkg{tomotopy})}

The baseline: plain LDA on the same 2000 posts, independent collapsed-Gibbs samplers (Java \pkg{MALLET}, \pkg{tomotopy}, and \pkg{topica}). Aligned topic-word cosine vs \pkg{topica} --- Java MALLET 0.822, tomotopy 0.803.

\begin{table}[ht]\centering\footnotesize
\caption{LDA topics on poliblog (K=10); all collapsed Gibbs samplers.}
\label{tab:app:lda}
\begin{tabular}{@{}r >{\raggedright\arraybackslash}p{0.29\textwidth} >{\raggedright\arraybackslash}p{0.29\textwidth} >{\raggedright\arraybackslash}p{0.29\textwidth}@{}}
\toprule
 & Java \pkg{MALLET} & \pkg{tomotopy} & \pkg{topica} \\
\midrule
1 & obama, mccain, campaign, john, barack, palin & obama, campaign, barack, clinton, hillari, polit & obama, mccain, campaign, john, barack, palin \\
2 & iraq, war, militari, forc, iran, bush & iraq, war, militari, iran, forc, bush & iraq, war, militari, forc, will, iran \\
3 & obama, democrat, vote, republican, state, voter & obama, vote, voter, poll, state, democrat & obama, hillari, democrat, vote, elect, state \\
4 & tax, govern, money, will, american, plan & will, tax, year, american, govern, economi & will, tax, american, govern, econom, economi \\
5 & report, time, news, new, stori, york & report, time, stori, media, news, new & report, news, stori, time, media, said \\
6 & like, think, just, know, say, get & will, like, make, can, one, just & get, like, can, just, want, one \\
7 & bush, hous, presid, senat, administr, state & bush, administr, state, law, presid, court & bush, presid, administr, state, law, offic \\
8 & peopl, american, women, black, polit, school & american, peopl, countri, america, year, world & american, america, will, countri, peopl, world \\
9 & will, one, make, polit, even, much & mccain, palin, john, campaign, presid, bush & polit, right, peopl, one, like, even \\
10 & will, can, american, peopl, work, countri & senat, democrat, vote, republican, hous, bill & senat, democrat, republican, bill, vote, support \\
\bottomrule
\end{tabular}
\end{table}


\subsection[Non-negative matrix factorization (vs scikit-learn)]{Non-negative matrix factorization (vs \pkg{scikit-learn})}

Non-negative factorization of the document-term matrix; \pkg{topica} reimplements \pkg{scikit-learn}'s multiplicative updates with the same NNDSVD start. Aligned topic-word cosine \textbf{0.980}.

\begin{table}[ht]\centering\footnotesize
\caption{NMF topics on poliblog (K=10); Frobenius loss, NNDSVD init.}
\label{tab:app:nmf}
\begin{tabular}{@{}r >{\raggedright\arraybackslash}p{0.40\textwidth} >{\raggedright\arraybackslash}p{0.40\textwidth}@{}}
\toprule
 & \pkg{scikit-learn} & \pkg{topica} \\
\midrule
1 & pentagon, build, dead, confirm, armi, virginia, navi, civilian & pentagon, build, dead, confirm, armi, virginia, navi, civilian \\
2 & mccain, campaign, john, romney, talk, debat, point, ask & mccain, campaign, john, romney, talk, debat, point, tax \\
3 & obama, barack, campaign, polici, wright, polit, senat, poll & obama, barack, campaign, polici, poll, wright, point, polit \\
4 & democrat, vote, elect, voter, state, republican, parti, primari & democrat, vote, elect, voter, state, republican, parti, poll \\
5 & will, one, like, presid, plan, make, end, new & will, one, like, new, also, end, plan, may \\
6 & like, one, think, say, peopl, get, stori, know & peopl, one, like, think, say, know, just, get \\
7 & hillari, clinton, obama, eastern, edward, updat, say, bill & hillari, clinton, eastern, edward, updat, say, talk, presid \\
8 & american, tax, economi, can, year, make, job, need & american, will, tax, economi, can, job, make, year \\
9 & bush, iraq, presid, said, war, administr, report, govern & bush, presid, said, iraq, administr, report, hous, govern \\
10 & world, peopl, must, nation, america, american, war, state & world, state, must, nation, iraq, war, peopl, unit \\
\bottomrule
\end{tabular}
\end{table}


\subsection[Latent semantic analysis (vs scikit-learn)]{Latent semantic analysis (vs \pkg{scikit-learn})}

Truncated SVD of the tf-idf matrix; \pkg{topica} matches \pkg{scikit-learn}'s \code{TruncatedSVD} (signed loadings, top terms by absolute value). Aligned cosine \textbf{0.990}.

\begin{table}[ht]\centering\footnotesize
\caption{LSA components on poliblog (K=10); top terms by absolute loading.}
\label{tab:app:lsa}
\begin{tabular}{@{}r >{\raggedright\arraybackslash}p{0.40\textwidth} >{\raggedright\arraybackslash}p{0.40\textwidth}@{}}
\toprule
 & \pkg{scikit-learn} & \pkg{topica} \\
\midrule
1 & obama, mccain, will, campaign, democrat, one, said, bush & obama, mccain, will, campaign, democrat, one, said, bush \\
2 & obama, mccain, campaign, hillari, poll, bush, iraq, govern & obama, mccain, campaign, hillari, poll, bush, iraq, govern \\
3 & mccain, obama, hillari, clinton, palin, john, democrat, bush & mccain, obama, hillari, clinton, palin, john, democrat, bush \\
4 & obama, palin, iraq, iraqi, iran, vote, republican, war & obama, iraq, palin, iraqi, iran, vote, republican, war \\
5 & palin, vote, democrat, poll, tax, sarah, bush, obama & palin, vote, democrat, poll, tax, bush, sarah, obama \\
6 & obama, vote, tax, poll, iraq, palin, iraqi, democrat & obama, vote, tax, iraq, poll, palin, iraqi, democrat \\
7 & blagojevich, senat, bush, hous, hillari, administr, investig, presid & bush, senat, blagojevich, hous, presid, administr, hillari, palin \\
8 & palin, mccain, obama, oil, sarah, iraqi, percent, poll & palin, mccain, oil, iraqi, sarah, obama, iraq, alaska \\
9 & hillari, clinton, poll, iraq, bush, israel, palin, blagojevich & hillari, clinton, poll, bush, israel, iraq, blagojevich, vote \\
10 & israel, iraqi, iran, iraq, blagojevich, hama, bush, report & iran, iraqi, israel, iraq, nuclear, hillari, hama, presid \\
\bottomrule
\end{tabular}
\end{table}

\noindent\emph{Singular values} (the variance each component captures), first 5:
\begin{center}\footnotesize\begin{tabular}{@{}l r r r r r @{}}\toprule
 & 1 & 2 & 3 & 4 & 5 \\ \midrule
\pkg{scikit-learn} & 12.1 & 5.4 & 4.5 & 4.1 & 4.0 \\
\pkg{topica} & 12.1 & 5.4 & 4.5 & 4.1 & 4.0 \\
\bottomrule\end{tabular}\end{center}


\subsection[Structural topic model (vs R stm)]{Structural topic model (vs \proglang{R}~\pkg{stm})}

Both engines fit \code{prevalence = \textasciitilde{} rating + s(day)} on the same 2000 posts and 2632 word types. Aligned topic-word cosine \textbf{0.986}.

\begin{table}[ht]\centering\footnotesize
\caption{STM topics on poliblog (K=10), \proglang{R}~\pkg{stm} vs \pkg{topica}, aligned by topic-word cosine.}
\label{tab:app:stm}
\begin{tabular}{@{}r >{\raggedright\arraybackslash}p{0.40\textwidth} >{\raggedright\arraybackslash}p{0.40\textwidth}@{}}
\toprule
 & \proglang{R}~\pkg{stm} & \pkg{topica} \\
\midrule
1 & will, tax, american, year, economi, govern, econom, can & will, tax, american, year, economi, govern, econom, can \\
2 & obama, barack, campaign, polit, report, senat, wright, said & obama, barack, campaign, polit, report, senat, wright, chicago \\
3 & mccain, senat, john, campaign, said, republican, democrat, bill & mccain, senat, john, said, campaign, republican, bill, democrat \\
4 & one, like, get, time, just, know, think, peopl & one, like, get, time, just, peopl, think, know \\
5 & palin, peopl, right, american, one, will, think, romney & palin, peopl, right, american, will, one, romney, school \\
6 & bush, administr, presid, said, court, law, hous, report & bush, presid, administr, said, court, hous, law, report \\
7 & obama, hillari, mccain, democrat, poll, will, clinton, voter & obama, hillari, mccain, democrat, clinton, poll, will, campaign \\
8 & vote, elect, state, democrat, will, immigr, voter, republican & vote, elect, state, democrat, will, immigr, voter, republican \\
9 & will, world, iran, israel, attack, terrorist, nation, nuclear & will, world, iran, israel, attack, build, terrorist, nuclear \\
10 & iraq, war, militari, iraqi, troop, report, forc, armi & iraq, war, forc, iraqi, troop, militari, bush, secur \\
\bottomrule
\end{tabular}
\end{table}

\noindent\emph{Prevalence by ideology.} Mean topic share, Liberal minus Conservative posts (positive = more Liberal):
\begin{center}\footnotesize\begin{tabular}{@{}p{0.5\textwidth} r r@{}}
\toprule
topic (\pkg{topica} top words) & \pkg{stm} & \pkg{topica} \\
\midrule
mccain, senat, john, said & +0.098 & +0.096 \\
will, world, iran, israel & -0.088 & -0.078 \\
one, like, get, time & -0.068 & -0.069 \\
\bottomrule
\end{tabular}\end{center}


\subsection[Correlated topic model (vs R stm, no covariates)]{Correlated topic model (vs \proglang{R}~\pkg{stm}, no covariates)}

No covariates, so this is the logistic-normal (correlated) topic model. Aligned cosine vs \pkg{topica} --- R stm 0.986.

\begin{table}[ht]\centering\footnotesize
\caption{CTM topics on poliblog (K=10); the logistic-normal (correlated) topic model.}
\label{tab:app:ctm}
\begin{tabular}{@{}r >{\raggedright\arraybackslash}p{0.44\textwidth} >{\raggedright\arraybackslash}p{0.44\textwidth}@{}}
\toprule
 & \proglang{R}~\pkg{stm} & \pkg{topica} \\
\midrule
1 & will, tax, american, year, govern, economi, econom, can & will, tax, american, year, economi, govern, econom, can \\
2 & mccain, senat, john, republican, said, campaign, democrat, bill & mccain, senat, john, republican, said, campaign, democrat, bill \\
3 & one, like, get, time, just, know, think, peopl & one, like, get, just, time, peopl, know, think \\
4 & obama, campaign, barack, polit, clinton, report, said, senat & obama, campaign, barack, polit, report, senat, said, wright \\
5 & bush, administr, presid, said, court, law, hous, report & bush, administr, presid, said, court, hous, law, report \\
6 & palin, peopl, right, american, one, will, romney, think & palin, peopl, right, will, american, one, romney, school \\
7 & obama, hillari, mccain, poll, democrat, will, voter, clinton & obama, hillari, mccain, poll, democrat, clinton, will, voter \\
8 & vote, elect, state, democrat, will, immigr, voter, republican & elect, vote, state, democrat, will, immigr, voter, republican \\
9 & will, world, iran, israel, attack, terrorist, nation, nuclear & world, will, iran, israel, attack, build, terrorist, nuclear \\
10 & iraq, war, militari, iraqi, troop, report, forc, armi & iraq, war, forc, iraqi, troop, militari, bush, secur \\
\bottomrule
\end{tabular}
\end{table}

\noindent\emph{Topic correlations} (the structure CTM adds over LDA), strongest pairs in \pkg{topica}:
\begin{center}\footnotesize\begin{tabular}{@{}>{\raggedright\arraybackslash}p{0.7\textwidth} r@{}}
\toprule topic pair & corr \\ \midrule
one, like, get / mccain, senat, john & -0.24 \\
obama, hillari, mccain / bush, administr, presid & -0.23 \\
world, will, iran / mccain, senat, john & -0.23 \\
\bottomrule\end{tabular}\end{center}


\subsection[Content covariate / SAGE (vs R stm)]{Content covariate / SAGE (vs \proglang{R}~\pkg{stm})}

\code{content = \textasciitilde{} rating}: each topic's words shift by ideology. Marginal aligned cosine \textbf{0.964}.

\begin{table}[ht]\centering\footnotesize
\caption{Content-model (SAGE) topics on poliblog (K=10); marginal topic-word averaged over rating levels.}
\label{tab:app:content}
\begin{tabular}{@{}r >{\raggedright\arraybackslash}p{0.40\textwidth} >{\raggedright\arraybackslash}p{0.40\textwidth}@{}}
\toprule
 & \proglang{R}~\pkg{stm} & \pkg{topica} \\
\midrule
1 & will, tax, year, american, govern, economi, econom, can & will, tax, american, year, govern, economi, can, make \\
2 & obama, campaign, barack, hillari, said, polit, clinton, say & obama, campaign, barack, polit, hillari, said, clinton, say \\
3 & one, like, think, just, get, peopl, know, time & one, like, get, time, think, peopl, just, know \\
4 & elect, vote, state, democrat, will, voter, republican, law & vote, elect, state, democrat, will, voter, republican, one \\
5 & mccain, obama, campaign, john, poll, will, state, hillari & obama, mccain, campaign, john, will, poll, state, hillari \\
6 & bush, presid, administr, iran, said, will, nation, state & bush, presid, administr, will, said, hous, iran, nation \\
7 & palin, mccain, romney, said, sarah, governor, right, say & palin, said, mccain, say, peopl, romney, right, one \\
8 & will, world, peopl, american, countri, israel, war, america & will, world, war, peopl, american, militari, one, israel \\
9 & senat, democrat, republican, vote, bill, hous, sen, parti & senat, democrat, vote, republican, obama, bill, mccain, will \\
10 & iraq, war, iraqi, troop, build, report, pentagon, armi & iraq, iraqi, war, troop, forc, secur, said, will \\
\bottomrule
\end{tabular}
\end{table}

\noindent\emph{Per-group wording} of one topic in \pkg{topica} (the SAGE deviation, TV distance 0.26):
\begin{center}\footnotesize\begin{tabular}{@{}l p{0.7\textwidth}@{}}\toprule
Conservative & obama, mccain, will, democrat, campaign, hillari, clinton, poll \\
Liberal & mccain, obama, campaign, john, poll, new, state, will \\
\bottomrule\end{tabular}\end{center}


\subsection[Dirichlet-multinomial regression (vs tomotopy)]{Dirichlet-multinomial regression (vs \pkg{tomotopy})}

The metadata covariate is the post's ideology (\code{rating}); DMR makes each document's topic prior log-linear in it. Aligned cosine \textbf{0.817} --- at the ceiling for independent DMR samplers: two \pkg{tomotopy} runs from different seeds agree only 0.714 with each other (these are random-initialized collapsed-Gibbs fits, so the dominant topics coincide while the diffuse tail differs run to run).

\begin{table}[ht]\centering\footnotesize
\caption{DMR topics on poliblog (K=10), \pkg{tomotopy} vs \pkg{topica} (both Gibbs with a metadata-driven document prior).}
\label{tab:app:dmr}
\begin{tabular}{@{}r >{\raggedright\arraybackslash}p{0.40\textwidth} >{\raggedright\arraybackslash}p{0.40\textwidth}@{}}
\toprule
 & \pkg{tomotopy} & \pkg{topica} \\
\midrule
1 & iraq, war, militari, iran, forc, bush, said, troop & iraq, war, militari, iran, forc, attack, iraqi, troop \\
2 & obama, vote, state, poll, voter, democrat, elect, mccain & obama, vote, state, democrat, voter, poll, elect, will \\
3 & will, tax, economi, govern, american, year, plan, econom & tax, will, american, govern, economi, econom, plan, year \\
4 & bush, state, presid, administr, law, court, investig, hous & bush, presid, administr, state, law, court, hous, investig \\
5 & will, like, one, just, can, make, get, peopl & will, can, get, one, like, just, know, make \\
6 & report, time, media, stori, new, news, one, day & report, said, news, time, stori, new, media, one \\
7 & senat, democrat, vote, republican, hous, bill, committe, blagojevich & senat, democrat, republican, vote, bill, hous, parti, romney \\
8 & obama, campaign, barack, clinton, hillari, polit, said, senat & obama, mccain, campaign, john, barack, palin, said, hillari \\
9 & american, peopl, countri, america, world, will, nation, year & american, america, world, countri, build, school, peopl, famili \\
10 & mccain, palin, john, campaign, said, say, bush, presid & polit, right, one, peopl, mani, believ, issu, like \\
\bottomrule
\end{tabular}
\end{table}

\noindent\emph{Prevalence by ideology.} Mean topic share, Liberal minus Conservative posts (positive = more Liberal):
\begin{center}\footnotesize\begin{tabular}{@{}p{0.5\textwidth} r r@{}}
\toprule
topic (\pkg{topica} top words) & \pkg{tomotopy} & \pkg{topica} \\
\midrule
will, can, get, one & -0.061 & -0.061 \\
polit, right, one, peopl & +0.048 & -0.010 \\
report, said, news, time & -0.026 & +0.008 \\
\bottomrule
\end{tabular}\end{center}


\subsection[Generalized DMR / continuous metadata (vs tomotopy)]{Generalized DMR / continuous metadata (vs \pkg{tomotopy})}

DMR generalized to \emph{continuous} metadata (post day, degree-3 Legendre basis). Topic-word aligned: cosine \textbf{0.818}.

\begin{table}[ht]\centering\footnotesize
\caption{g-DMR topics on poliblog (K=10); continuous-time (post day) metadata via a Legendre basis.}
\label{tab:app:gdmr}
\begin{tabular}{@{}r >{\raggedright\arraybackslash}p{0.40\textwidth} >{\raggedright\arraybackslash}p{0.40\textwidth}@{}}
\toprule
 & \pkg{tomotopy} & \pkg{topica} \\
\midrule
1 & iraq, war, militari, forc, iran, attack, bush, secur & iraq, war, militari, iran, forc, bush, iraqi, troop \\
2 & will, tax, year, american, economi, econom, govern, plan & tax, will, american, govern, economi, econom, plan, year \\
3 & obama, vote, democrat, poll, state, voter, elect, mccain & obama, vote, state, democrat, voter, poll, elect, hillari \\
4 & bush, state, presid, administr, law, court, investig, depart & bush, presid, administr, state, law, court, hous, investig \\
5 & like, will, one, just, can, make, get, peopl & will, can, get, one, like, just, want, make \\
6 & report, time, new, stori, one, news, media, show & report, said, new, time, news, stori, media, press \\
7 & senat, vote, democrat, republican, hous, bill, committe, sen & senat, democrat, republican, vote, bill, hous, support, romney \\
8 & american, peopl, america, countri, will, world, nation, right & american, america, countri, world, nation, peopl, build, school \\
9 & obama, campaign, barack, hillari, clinton, polit, said, say & obama, mccain, campaign, john, barack, palin, said, presid \\
10 & mccain, palin, john, campaign, said, say, bush, romney & polit, right, one, mani, peopl, support, think, like \\
\bottomrule
\end{tabular}
\end{table}

\noindent\emph{Prevalence over time.} Topic share at late minus early 2008 (positive = rising through the year):
\begin{center}\footnotesize\begin{tabular}{@{}p{0.5\textwidth} r r@{}}\toprule
topic (\pkg{topica} top words) & \pkg{tomotopy} & \pkg{topica} \\ \midrule
iraq, war, militari, iran & -0.111 & -0.034 \\
will, can, get, one & +0.007 & -0.033 \\
tax, will, american, govern & +0.052 & +0.032 \\
\bottomrule\end{tabular}\end{center}


\subsection[Keyword-assisted topic model (vs R keyATM)]{Keyword-assisted topic model (vs \proglang{R}~\pkg{keyATM})}

4 seeded keyword topics plus 6 free topics on 2000 posts. The anchored topics line up by construction; their recovered words agree at per-topic cosine \textbf{0.889} (all 10 topics, best-aligned: 0.661).

\begin{table}[ht]\centering\footnotesize
\caption{keyATM keyword topics on poliblog, \proglang{R}~\pkg{keyATM} vs \pkg{topica}; each topic anchored to its seed words (index-aligned).}
\label{tab:app:keyatm}
\begin{tabular}{@{}>{\raggedright\arraybackslash}p{0.18\textwidth} >{\raggedright\arraybackslash}p{0.36\textwidth} >{\raggedright\arraybackslash}p{0.36\textwidth}@{}}
\toprule
seed keywords & \proglang{R}~\pkg{keyATM} & \pkg{topica} \\
\midrule
tax, economi, econom, market, spend, budget & tax, will, econom, economi, govern, plan, american, year & tax, economi, will, econom, billion, govern, oil, plan \\
iraq, iraqi, war, troop, militari, surg & iraq, war, militari, iran, forc, iraqi, bush, secur & iraq, war, iran, militari, bush, iraqi, forc, said \\
obama, mccain, vote, voter, campaign, elect & obama, mccain, campaign, democrat, vote, will, republican, john & obama, mccain, campaign, democrat, vote, hillari, republican, barack \\
abort, gay, marriag, religi, church, famili & obama, famili, wright, women, school, black, church, america & obama, wright, blagojevich, black, church, women, abort, said \\
\bottomrule
\end{tabular}
\end{table}


\subsection[Supervised LDA (vs tomotopy)]{Supervised LDA (vs \pkg{tomotopy})}

LDA with a response regression: topics are shaped to predict ideology. Topic-word aligned: cosine \textbf{0.727}.

\begin{table}[ht]\centering\footnotesize
\caption{sLDA topics on poliblog (K=10); response = ideology (Liberal\,=\,1).}
\label{tab:app:slda}
\begin{tabular}{@{}r >{\raggedright\arraybackslash}p{0.40\textwidth} >{\raggedright\arraybackslash}p{0.40\textwidth}@{}}
\toprule
 & \pkg{tomotopy} & \pkg{topica} \\
\midrule
1 & mccain, palin, john, campaign, said, debat, sen, presid & mccain, john, campaign, bush, said, presid, sen, say \\
2 & obama, campaign, barack, hillari, democrat, clinton, voter, state & obama, barack, campaign, senat, polit, clinton, will, hillari \\
3 & iraq, war, militari, iran, bush, forc, said, attack & iraq, war, will, militari, iran, forc, iraqi, govern \\
4 & state, administr, presid, bush, law, said, court, hous & law, report, court, administr, state, offic, hous, investig \\
5 & tax, will, year, american, economi, govern, econom, plan & will, oil, hous, energi, money, compani, price, financi \\
6 & think, say, know, like, just, peopl, get, said & one, get, like, time, stori, media, just, question \\
7 & will, one, can, make, even, like, may, much & peopl, will, right, think, like, polit, one, say \\
8 & american, peopl, america, world, countri, nation, school, year & will, tax, american, can, peopl, work, econom, year \\
9 & vote, democrat, republican, senat, hous, bill, gop, parti & obama, vote, democrat, will, voter, poll, elect, hillari \\
10 & report, time, news, new, media, stori, one, day & palin, romney, one, women, school, year, sarah, governor \\
\bottomrule
\end{tabular}
\end{table}

\noindent\emph{Response coefficients.} How each topic moves the ideology response (sign = direction; the two engines scale the latent differently, so compare sign and ranking):
\begin{center}\footnotesize\begin{tabular}{@{}p{0.5\textwidth} r r@{}}\toprule
topic (\pkg{topica} top words) & \pkg{tomotopy} & \pkg{topica} \\ \midrule
mccain, john, campaign, bush & +1.68 & +2.00 \\
peopl, will, right, think & -0.61 & +0.84 \\
law, report, court, administr & +1.17 & +0.80 \\
\bottomrule\end{tabular}\end{center}


\subsection[Labeled LDA (vs tomotopy)]{Labeled LDA (vs \pkg{tomotopy})}

Supervised topics pinned to document labels: the topic set \emph{is} the label set (here ideology), so the topics are index-aligned by construction. Per-label cosine \textbf{1.000}.

\begin{table}[ht]\centering\footnotesize
\caption{Labeled-LDA topics on poliblog: one topic per ideology label, \pkg{tomotopy} vs \pkg{topica} (index-aligned by label).}
\label{tab:app:labeledlda}
\begin{tabular}{@{}l >{\raggedright\arraybackslash}p{0.36\textwidth} >{\raggedright\arraybackslash}p{0.36\textwidth}@{}}
\toprule
label & \pkg{tomotopy} & \pkg{topica} \\
\midrule
Conservative & obama, will, mccain, one, can, time, like, peopl & obama, will, mccain, one, can, time, like, peopl \\
Liberal & mccain, obama, will, said, campaign, one, say, presid & mccain, obama, will, said, campaign, one, say, presid \\
\bottomrule
\end{tabular}
\end{table}


\subsection[Hierarchical Dirichlet process (vs tomotopy)]{Hierarchical Dirichlet process (vs \pkg{tomotopy})}

The nonparametric model: both engines \emph{infer} the topic count rather than fixing it, and they land in the same range --- \pkg{tomotopy} discovered \textbf{37} live topics, \pkg{topica} \textbf{28}, a quantity neither was told. That count recovery is HDP's headline check. Topic identities, by contrast, are seed-sensitive in any HDP sampler: \pkg{topica}'s 10 most prevalent topics match \pkg{tomotopy}'s at cosine 0.606, which is its own seed-to-seed ceiling (0.585 between two \pkg{topica} seeds) --- the cross-engine match is as tight as the model is reproducible with itself.

\begin{table}[ht]\centering\footnotesize
\caption{HDP topics on poliblog: \pkg{topica}'s 10 most prevalent topics beside their best-matched \pkg{tomotopy} topic.}
\label{tab:app:hdp}
\begin{tabular}{@{}r >{\raggedright\arraybackslash}p{0.40\textwidth} >{\raggedright\arraybackslash}p{0.40\textwidth}@{}}
\toprule
 & \pkg{tomotopy} & \pkg{topica} \\
\midrule
1 & iraq, will, bush, presid, war, obama, said, nation & will, can, one, like, think, even, want, say \\
2 & one, will, right, peopl, like, say, can, just & peopl, make, said, get, just, right, support, presid \\
3 & will, iraq, war, said, american, one, militari, attack & iraq, war, militari, iran, bush, forc, govern, secur \\
4 & obama, polit, one, time, like, campaign, just, stori & time, new, stori, live, man, post, american, day \\
5 & obama, mccain, campaign, will, hillari, democrat, poll, vote & obama, vote, democrat, voter, poll, state, elect, will \\
6 & mccain, palin, obama, peopl, think, one, say, know & obama, mccain, campaign, john, barack, palin, senat, say \\
7 & report, state, senat, pentagon, hous, investig, build, confirm & state, report, court, law, investig, administr, offic, feder \\
8 & will, tax, american, mccain, year, can, govern, make & tax, will, govern, econom, economi, american, plan, billion \\
9 & mccain, obama, democrat, senat, will, said, campaign, vote & senat, democrat, vote, republican, hous, bill, bush, presid \\
10 & oil, water, product, reid, find, energi, nuclear, system & oil, energi, price, compani, drill, global, state, increas \\
\bottomrule
\end{tabular}
\end{table}


\subsection[Pachinko allocation (vs tomotopy)]{Pachinko allocation (vs \pkg{tomotopy})}

A DAG of 3 super-topics over 10 shared sub-topics, capturing topic correlation through the hierarchy. Sub-topics aligned: cosine \textbf{0.755}.

\begin{table}[ht]\centering\footnotesize
\caption{Pachinko sub-topics on poliblog (3 super-topics over 10 sub-topics), \pkg{tomotopy} vs \pkg{topica}.}
\label{tab:app:pa}
\begin{tabular}{@{}r >{\raggedright\arraybackslash}p{0.40\textwidth} >{\raggedright\arraybackslash}p{0.40\textwidth}@{}}
\toprule
 & \pkg{tomotopy} & \pkg{topica} \\
\midrule
1 & tax, will, american, economi, year, econom, plan, govern & tax, american, will, govern, economi, econom, plan, year \\
2 & iraq, war, militari, iran, forc, attack, bush, iraqi & iraq, war, militari, forc, iran, american, attack, bush \\
3 & obama, democrat, vote, voter, poll, state, elect, will & democrat, obama, vote, state, republican, voter, elect, poll \\
4 & get, one, can, like, just, even, make, will & one, like, peopl, say, just, think, make, can \\
5 & obama, mccain, campaign, barack, john, said, hillari, senat & obama, campaign, barack, hillari, clinton, polit, candid, question \\
6 & report, palin, media, news, time, stori, new, press & report, stori, news, time, new, media, york, show \\
7 & peopl, american, countri, america, right, believ, polit, one & american, peopl, america, world, countri, school, will, build \\
8 & court, law, offic, investig, state, report, administr, depart & hous, bush, presid, state, administr, law, senat, said \\
9 & will, nation, new, state, includ, power, also, govern & will, can, year, time, nation, take, even, also \\
10 & bush, senat, republican, democrat, hous, presid, vote, bill & mccain, palin, john, campaign, said, sen, senat, today \\
\bottomrule
\end{tabular}
\end{table}

\noindent\emph{Super-topic structure} in \pkg{topica}: each super-topic's two strongest sub-topics (the Pachinko DAG that plain LDA lacks):
\begin{center}\footnotesize\begin{tabular}{@{}l p{0.72\textwidth}@{}}\toprule
super-topic & top sub-topics (by association) \\ \midrule
1 & one, like, peopl; american, peopl, america \\
2 & one, like, peopl; will, can, year \\
3 & one, like, peopl; iraq, war, militari \\
\bottomrule\end{tabular}\end{center}


\subsection[Dynamic topic model (vs tomotopy)]{Dynamic topic model (vs \pkg{tomotopy})}

Topics whose word distributions drift across 4 time slices of 2008. At the final slice the aligned topic-word cosine is \textbf{0.735}: with a converged reference the dominant themes line up, though \pkg{tomotopy}'s \code{DTModel} separates topics less sharply than \pkg{topica} on this corpus, which caps the cosine. The model's defining behaviour --- how a topic's vocabulary shifts over the year --- is shown below.

\begin{table}[ht]\centering\footnotesize
\caption{DTM topics at the final slice on poliblog (K=10, 4 time slices of 2008).}
\label{tab:app:dtm}
\begin{tabular}{@{}r >{\raggedright\arraybackslash}p{0.40\textwidth} >{\raggedright\arraybackslash}p{0.40\textwidth}@{}}
\toprule
 & \pkg{tomotopy} & \pkg{topica} \\
\midrule
1 & obama, state, say, one, said, barack, senat, campaign & obama, barack, campaign, senat, polit, will, clinton, one \\
2 & mccain, campaign, time, barack, one, john, republican, vote & mccain, john, campaign, said, bush, obama, presid, sen \\
3 & will, peopl, say, can, like, one, now, said & peopl, will, like, right, think, one, polit, say \\
4 & one, like, democrat, will, get, make, new, said & one, get, time, media, stori, like, just, news \\
5 & palin, peopl, like, year, can, time, one, just & palin, sarah, romney, women, one, school, governor, huckabe \\
6 & one, said, year, think, want, right, can, will & will, tax, american, can, peopl, year, work, economi \\
7 & report, vote, nation, one, said, american, new, peopl & law, court, administr, bush, report, hous, state, said \\
8 & democrat, elect, like, can, just, year, one, obama & obama, vote, democrat, will, poll, voter, elect, state \\
9 & senat, time, american, new, know, vote, democrat, year & iraq, war, will, iran, militari, forc, iraqi, govern \\
10 & will, can, new, make, presid, time, now, one & will, oil, hous, energi, compani, price, govern, year \\
\bottomrule
\end{tabular}
\end{table}

\noindent\emph{Word drift} within one \pkg{topica} topic across 2008 (first to last slice):
\begin{center}\footnotesize\begin{tabular}{@{}l p{0.72\textwidth}@{}}\toprule
rising & palin, sarah, alaska, women, research, governor \\
falling & romney, huckabe, fred, thompson, paul, global \\
\bottomrule\end{tabular}\end{center}


\subsection[Structural topic and sentiment-discourse (vs the authors' R reference)]{Structural topic and sentiment-discourse (vs the authors' \proglang{R} reference)}

The authors' published poliblog\,2008 fit vs \pkg{topica}, both read at the mean sentiment (the reference's \code{print.topWords} view). Aligned cosine \textbf{0.930}.

\begin{table}[ht]\centering\footnotesize
\caption{STS topics on the published poliblog fit (K=5), read at each topic's mean sentiment-discourse.}
\label{tab:app:sts}
\begin{tabular}{@{}r >{\raggedright\arraybackslash}p{0.40\textwidth} >{\raggedright\arraybackslash}p{0.40\textwidth}@{}}
\toprule
 & \proglang{R} reference & \pkg{topica} \\
\midrule
1 & obama, democrat, vote, will, republican, elect, hillari, clinton & obama, democrat, vote, will, hillari, elect, clinton, republican \\
2 & obama, mccain, campaign, said, barack, john, palin, say & obama, mccain, campaign, barack, palin, john, said, say \\
3 & one, peopl, like, will, can, just, right, think & will, one, peopl, like, can, time, year, american \\
4 & iraq, war, will, militari, bush, iran, american, said & bush, iraq, will, presid, said, govern, war, state \\
5 & will, bush, govern, hous, tax, year, american, presid & oil, energi, year, democrat, will, price, tax, hous \\
\bottomrule
\end{tabular}
\end{table}

\noindent\emph{Sentiment by ideology.} Effect of \code{rating} (Liberal) on each topic's sentiment-discourse latent in \pkg{topica} (positive = more positive among Liberal posts) — the dimension STS adds over STM:
\begin{center}\footnotesize\begin{tabular}{@{}p{0.6\textwidth} r@{}}\toprule
topic (top words) & sentiment effect \\ \midrule
obama, mccain, campaign, barack, palin & +0.429 \\
obama, democrat, vote, republican, will & +0.195 \\
allah, qur, sura, ibn, muhammad & +0.072 \\
\bottomrule\end{tabular}\end{center}



\end{document}
